\section{Applications and admissions}
	\setcounter{table}{0}
	\renewcommand{\thetable}{F\arabic{table}}
	
	\setcounter{figure}{0}
	\renewcommand{\thefigure}{F\arabic{figure}}
\label{application_lists}
\paragraph{Comparing regular applications and admissions between students in treated and control schools.} PACE had null effects on the proportion of students sending a regular application and receiving a regular admission (Table \ref{tab: TEapplicationsadmissionsregular}). Descriptive statistics show that regular-channel applicants from treated schools tend to apply and be admitted to more selective majors than regular-channel applicants from control schools (first three columns of Table \ref{tab:summaryapplications} and Figure \ref{fig:histogram_quality_regular_T_C_comb}), but the differences in the application and admission patterns are close to zero and statistically insignificant once we control for the different pool of applicants across treatment groups by including students' baseline characteristics (Panels A and C, columns (3) and (6) of Table \ref{tab: TEapplicationsassignmentstop15}).\footnote{There is a statistically significant difference at the 10\% level for applications in the all-student sample, but the significance disappears once we account for multiple hypothesis testing (RW-adjusted p-value: 0.405). } Descriptive statistics also suggest that regular applicants from treated schools apply and are admitted to programs that are closer to their high school on average (first three columns of Table \ref{tab:summaryapplications}), because they apply at lower rates to programs that are more than 500km away (Figure \ref{fig:histogram_distance_regular_T_C_comb}). But these differences are statistically insignificant (columns (1) and (4) of Table \ref{tab: TEapplicationsassignmentstop15}). Finally, application and admission patterns across various majors are similar between treatment groups. The main exception is that students from treated schools are admitted at higher rates to natural sciences and lower rates to engineering---both STEM majors---compared to those from control schools (Figure \ref{fig:histogram_top13_application_regular_T_C_comb}). However, the lower regular-channel engineering admissions for treated school students are offset by higher PACE-channel engineering admissions (Figure \ref{fig:histogram_top13_application_T_PACE_regular_top15_comb}). And importantly, regular applicants from both groups apply and are admitted to STEM and non-STEM majors at nearly identical rates (first three columns of Table \ref{tab:summaryapplications}, columns (2) and (5) of Table \ref{tab: TEapplicationsassignmentstop15}). Therefore, we do not find significant differences in admission patterns through the regular channel across treatment groups.\par 
            

            


      
      \paragraph{Comparing applications and admissions across the regular and PACE channels for students in treated schools.} Top-performing students within their high school, who are those where PACE applications and admissions are concentrated, apply to more selective programs through the PACE than through the regular channel, but they are admitted to programs that are similarly selective across channels (Figure \ref{fig:histogram_quality_T_PACE_regular_top15_comb}, Panel A of Table \ref{tab:summaryapplications}, and column (3) of Table \ref{tab: diffapplicationsassignmentstreatedtop15}). They apply and are admitted to programs that are similarly distant from their high school across the regular and PACE channels (Figure \ref{fig:histogram_distance_top15_comb} and Panel A of Table \ref{tab:summaryapplications}), with only small and insignificant differences across channels (column (1) of Table \ref{tab: diffapplicationsassignmentstreatedtop15}). Finally, these students are slightly more likely to send applications to STEM programs through the PACE channel, although the difference is not statistically significant for the top choice once we account for multiple hypotheses testing (Panel A of Table \ref{tab:summaryapplications} and Panels A and B, column (2) of Table \ref{tab: diffapplicationsassignmentstreatedtop15}). This is driven by listing slightly more engineering and health programs and slightly fewer education and arts programs (Figure \ref{fig:histogram_top13_application_T_PACE_regular_top15_comb}). While there are some differences in the majors to which students are admitted across channels---notably, students are more likely to be admitted to engineering and less likely to be admitted to health programs through the PACE channel (Figure \ref{fig:histogram_top13_application_T_PACE_regular_top15_comb})---there are no statistically significant differences across channels in the STEM composition of the programs to which they are admitted (Panel C, column (2) of Table \ref{tab: diffapplicationsassignmentstreatedtop15}). Therefore, we do not find significant differences in the admission patterns through the regular and PACE channels for students in PACE schools.\par 


\subsection{Figures}

%% Selectivity 
\subsubsection{Selectivity}
\begin{figure}[H] %was [ht]
	\centering
\includegraphics[width=\linewidth]{output/figures/histogram_quality_regular_T_C_comb.png}
\caption{\label{fig:histogram_quality_regular_T_C_comb} Selectivity of programs to which students apply through the regular channel and where they are admitted across treatment and control groups. The left panel shows all students regardless of their ranking in their high school, the right panel shows the students who were in the top 15\% of their high school GPA ranking at baseline.}
\end{figure}

\begin{figure}[H] %was [ht]
	\centering
\includegraphics[width=\linewidth]{output/figures/histogram_quality_T_PACE_regular_top15_comb.png}
\caption{\label{fig:histogram_quality_T_PACE_regular_top15_comb} Selectivity of programs to which top 15\% treated students apply and where they are admitted across regular and PACE application channels. These students attended treated schools and were in the top 15\% of their high school GPA ranking at baseline.}
\end{figure}


%%% Location 
\subsubsection{Location }
\begin{figure}[H] %was [ht]
	\centering
\includegraphics[width=\linewidth]{output/figures/histogram_distance_regular_T_C_comb.png}
\caption{\label{fig:histogram_distance_regular_T_C_comb} Location of programs to which students apply through the regular channel and where they are admitted through the regular channel across treatment and control groups. The left panel shows all students regardless of their ranking in their high school, the right panel shows the students who were in the top 15\% of their high school GPA ranking at baseline.}
\end{figure}

\begin{figure}[H] %was [ht]
	\centering
\includegraphics[width=\linewidth]{output/figures/histogram_distance_top15_comb.png}
\caption{\label{fig:histogram_distance_top15_comb} Location of programs to which top 15\% treated students apply and where they are admitted across regular and PACE application channels. These students attended treated schools and were in the top 15\% of their high school GPA ranking at baseline.}
\end{figure}


%% Study field

\subsubsection{Study field}\label{sec:appendixSTEMfigures}
\begin{figure}[H] %was [ht]
	\centering
\includegraphics[width=\linewidth]{output/figures/histogram_field_regular_T_C_comb.png}
\caption{\label{fig:histogram_top13_application_regular_T_C_comb} Study field to which students apply through the regular channel and where they are admitted through the regular channel across treatment and control groups. The left panel shows all students regardless of their ranking in their high school, the right panel shows the students who were in the top 15\% of their high school GPA ranking at baseline. The average study field of listed programs is computed by taking the average across students of the fraction of programs listed by each student belonging to that study field.}
\end{figure}

\begin{figure}[H] %was [ht]
	\centering
\includegraphics[width=\linewidth]{output/figures/histogram_field_T_PACE_regular_top15_comb.png}
\caption{\label{fig:histogram_top13_application_T_PACE_regular_top15_comb} Study field to which top 15\% treated students apply and where they are admitted across regular and PACE application channels.  These students attended treated schools and were in the top 15\% of their high school GPA ranking at baseline. The average study field of listed programs is computed by taking the average across students of the fraction of programs listed by each student belonging to that study field.}
\end{figure}














\subsection{Tables}




\input{output/tables/TE_regular_admissions_enrollments_f}

\input{output/tables/summary_stats_applications_assignments_top15_all}


\input{output/tables/TE_applications_assignments_top15}

\input{output/tables/differences_pace_regular_for_treated_top15}




\pagebreak
\newpage



































